\chapter{Reading Routes by Level and Interest}

You need not work through all twenty chapters from the first page to the last. Mathematics is less a one-way street than a metro map, with several main lines and many interchanges. This appendix offers a default route and four routes by interest. Choose your own pace and preferences, and change lines whenever you like.

\section{The Default Route: Four Steps}

If you have no particular preference, these four steps provide a steady path, each preparing the next.

\textbf{Step one: Chapters $1$--$4$, foundations.} Start with algebraic symbols, equations, and functions to practise symbolic reasoning. These chapters supply the book's shared language, needed on any later route. Read them in order and consult Appendix B whenever elementary material is unfamiliar.

\textbf{Step two: Chapters $5$--$8$, going deeper.} Polynomials, transformations, and vectors put the first step's tools to deeper use. Here you first appreciate how the same object, such as a polynomial, can have several algebraic and geometric faces. Changing the face through which you view a problem is one of a mathematician's handiest moves.

\textbf{Step three: Chapters $9$--$12$, the language of higher mathematics.} Move from change to infinity, meeting intuitive pictures of limits, derivatives, and integrals. The keyword is approximation. Allow yourself time to become comfortable with rigorous statements about infinity. First grasp the graphs and numerical examples; rigor will follow. Letting intuition fall behind is the greater difficulty.

\textbf{Step four: Chapter $13$ onward, branching by interest.} From here, chapters no longer form a strict sequence. Number theory, topology, geometry, probability, and mathematical foundations develop their own lines. After step three, you can jump to whatever looks intriguing. The following routes are designed for that freedom.

\section{Four Routes by Interest}

\textbf{Route one: problem solving---Chapters $1$--$6$, $13$, and $14$.} If you enjoy the relief of solving a problem that initially seemed baffling, this route leads to mathematics' oldest hunting ground: integers. The first six chapters sharpen algebraic tools; Chapters $13$ and $14$ enter number theory, where a question may take one sentence to understand, such as a conjecture about primes, yet stump mathematicians for centuries. Keep scrap paper nearby and work each example yourself before reading on. The unit of problem-solving ability is not pages read, but sheets of working.

\textbf{Route two: shapes and spaces---Chapters $3$, $4$, $6$--$12$, $16$, and $17$.} If pictures feel friendlier than formulas, begin with coordinates and transformations, continue through function graphs, rates of change, and curves, and arrive at topology and differential geometry: geometry on rubber sheets and shortest paths in curved spaces. Draw often, even badly. Geometric intuition comes from making drawings, not merely looking at them.

\textbf{Route three: infinity and change---Chapters $8$--$12$, $15$, and $17$.} Infinity is one of mathematics' most captivating cliffs. What is an infinitesimal? How can infinitely many positive numbers have a finite sum? This route begins with rates of change, approaches limits and infinite series directly, then extends into related geometric and analytic ideas. Bring patience. For every conclusion about infinity, ask how it is approached step by step.

\textbf{Route four: codes, computation, and applications---Chapters $2$, $5$, $7$, $13$--$15$, $18$, and $19$.} If your question is ``what is mathematics useful for?'', begin with equations and algebraic structures, pass through number theory---a foundation of modern cryptography, with prime-number properties used in encrypted phone payments---and continue to probability and real-world modeling. Combine this route with Appendix D's experiments. Many ideas along it become immediately playable in a spreadsheet or a few lines of code.

\begin{lianjie}
The routes' connections around Chapter $13$ are no accident. Number theory uses elementary tools such as divisibility and remainders, yet connects to cryptography, geometry, and even analysis. Mathematics has many hidden passages between branches. The book's Connections boxes keep pointing them out. Wandering from one route onto another is a normal and rewarding part of reading, not something to correct.
\end{lianjie}

\section{Three Suggestions for Reading}

First, \textbf{getting stuck is normal}. The reflection questions provide hints only. Returning after several days of being stuck often teaches more than immediately reading an answer. Mathematicians themselves can spend years stuck on a problem. Second, \textbf{you need not understand everything at once}. Follow the main story on a first reading and save proofs for a second. Recognizing where the difficulty lies is progress. Third, \textbf{give hands-on work priority}. Stop for the One-Minute Experiments. Five minutes of doing often beats half an hour of staring.

The other four appendices are tools to keep beside you. Appendix A is a grammar of proof, useful when proofs first seem bewildering. Appendix B refreshes elementary tools such as fractions, equations, and coordinates. Appendix C is a symbol dictionary, and Appendix D provides experimental equipment. Together with this appendix, they form the book's support team, ready wherever your current chapter takes you.

\begin{pandeng}
To keep climbing after this book, choose the next mountain by route. Problem solving leads to competition mathematics and elementary number theory texts. Shapes lead to \emph{What Is Mathematics?} and introductory books on topology. Infinity leads to first courses in calculus and mathematical analysis. Applications lead to accessible books and open courses on cryptography and probability. With keywords in hand, libraries and the internet offer many paths.
\end{pandeng}

\section*{Reflection Questions}

\textbf{Observation}
\begin{sikao}
\item Browse the contents, choose the three chapters you most want to read, and identify their routes. Hint: a chapter can belong to several routes, and intersections are often especially interesting.
\end{sikao}

\textbf{Proof}
\begin{sikao}
\item This appendix has no theorem to prove, but try arguing to a classmate why Chapters $1$--$4$ precede almost every route. Hint: open any later chapter and count the tools it uses from those early chapters.
\end{sikao}

\textbf{Challenge}
\begin{sikao}
\item Design a route absent from this appendix, perhaps for someone interested in mathematical history and mathematicians' lives. List the chapter stops and explain why each deserves a visit. Hint: first decide what you want to see at the destination, then work backward to the preparation needed along the way.
\end{sikao}
